\documentclass[10pt,twocolumn]{article}

% ============================================================
% Packages
% ============================================================
\usepackage[utf8]{inputenc}
\usepackage[T1]{fontenc}
\usepackage{kotex}
\usepackage[margin=0.75in]{geometry}
\usepackage{amsmath,amssymb,amsfonts}
\usepackage{graphicx}
\usepackage{booktabs}
\usepackage{hyperref}
\usepackage{algorithm}
\usepackage{algorithmic}
\usepackage{xcolor}
\usepackage{enumitem}
\usepackage{caption}
\usepackage{tabularx}
\usepackage{multirow}
\usepackage{url}

\hypersetup{colorlinks=true,linkcolor=blue,citecolor=blue,urlcolor=blue}

\title{TurboQuant 실전 구현: 프로덕션 LLM 추론 엔진에서의\\3비트 KV 캐시 압축}

\author{
AmesianX\\
독립 연구자\\
\texttt{amesianx@gmail.com}\\
\url{https://powerhacker.net}\\
\url{https://github.com/AmesianX/TurboQuant}
}

\date{2026년 4월}

% ============================================================
\begin{document}
\maketitle

% ============================================================
\begin{abstract}
본 논문은 Walsh-Hadamard Transform(WHT)과 Lloyd-Max 양자화 기반 KV 캐시 압축 방법인 TurboQuant을, 가장 널리 배포된 오픈소스 LLM 추론 엔진 llama.cpp에 세계 최초로 프로덕션 수준으로 구현한 결과를 제시한다. 본 구현은 KV 캐시를 요소당 3비트로 압축하여 Gemma~4 26B 모델에서 FP16 대비 \textbf{4.2배 원시 메모리 절감}과 \textbf{6\% perplexity 저하}만을 달성했으며, 하이브리드 rank-value 압축을 통해 추론 시 유효 압축률이 \textbf{${\sim}4.74\times$}에 이른다. 원본 논문의 충실한 구현을 넘어 네 가지 독창적 기여를 한다: (1)~MMSE와 극치 통계학(EVT)에서 유도한 \emph{동적 어텐션 샤프닝} 공식 $\alpha(N) = 1 + c\sqrt{\ln N}$; (2)~flash attention 커널의 \emph{반정밀도 누적 버그} 발견 및 해결 — FP16 산술의 정밀도 손실이 IWHT 버터플라이에 의해 증폭되어 MoE expert routing이 비결정적으로 동작하는 문제; (3)~\emph{CUDA 네이티브 TriAttention} 통합 — Parseval 보존을 활용하여 양자화된 WHT 도메인 K 블록 위에서 직접 trig 기반 중요도 스코어를 계산하며, 참조 구현에서 지배적이던 GPU$\to$CPU 왕복을 제거하고 \emph{rank 축}(누구를 남길지)과 \emph{value 축}(원소당 비트 수)을 깔끔하게 분리; (4)~5개 head dimension(64--576)에 걸친 22개 양자화 타입 변형, double WHT per-head, 비대칭 K/V 구성, 그리고 TriAttention 스코어링에 필요한 trig 요약을 함께 담는 하이브리드 AMX3\_1 블록 레이아웃(WHT 50B + polar 58B = 108B)을 지원하는 완전한 엔지니어링 프레임워크. 전체 구현을 \url{https://github.com/AmesianX/TurboQuant}에 오픈소스로 공개한다.
\end{abstract}

% ============================================================
\section{서론}
\label{sec:intro}

긴 컨텍스트 윈도우를 가진 대규모 언어 모델(LLM)은 키-밸류(KV) 캐시에 상당한 GPU 메모리를 요구한다. 262K 컨텍스트의 26B 파라미터 Gemma~4 모델의 경우, FP16 KV 캐시만으로 5.4~GiB를 소비하며 이는 종종 모델 가중치 자체를 초과한다. Google DeepMind가 ICLR 2026에서 발표한 TurboQuant~\cite{turboquant2026}은 Walsh-Hadamard Transform(WHT) 후 Lloyd-Max 양자화를 적용하고, 선택적으로 QJL(Quantized Johnson-Lindenstrauss) 보정 항을 추가하여 KV 캐시 항목을 3--4비트로 압축한다.

이론적 프레임워크는 우아하지만, 본 연구 이전에 프로덕션 구현은 존재하지 않았다. 연구 논문과 실제 CUDA 커널 사이의 간극을 메우려면 논문이 다루지 않는 수많은 엔지니어링 문제를 해결해야 한다: WHT 버터플라이 연산을 위한 공유 메모리 레이아웃, 역변환의 정밀도 요구사항, 기존 flash attention 프레임워크와의 상호작용, 현대 아키텍처의 다양한 head dimension(64, 128, 256, 512, 576) 지원.

본 논문에서는 llama.cpp의 flash attention vector 커널(\texttt{fattn-vec}) 내 완전한 구현을 기술하고, 실제 모델에서의 성능을 보고하며, 원본 논문을 넘어서는 네 가지 기여를 제시한다:

\begin{enumerate}[leftmargin=*,topsep=2pt,itemsep=1pt]
\item \textbf{동적 어텐션 샤프닝} (\S\ref{sec:sharpening}): 양자화 노이즈가 softmax 분포를 평탄화한다. MMSE 이론과 EVT에서 컨텍스트 길이 적응형 보정 $\alpha(N) = 1 + c\sqrt{\ln N}$을 유도하여 경험적 튜닝을 제거했다.

\item \textbf{반정밀도 병리학} (\S\ref{sec:precision}): 업스트림 flash attention 커널의 FP16(\texttt{half2}) 누적 경로가 정밀도 손실을 야기하고, 이것이 IWHT 버터플라이에 의해 \emph{증폭}되어 MoE expert routing에서 치명적 비결정성을 초래함을 발견하고 해결했다. TBQ 밸류 타입에 FP32 누적을 강제하여 perplexity가 2\% 개선되었다.

\item \textbf{양자화된 K에서의 CUDA 네이티브 TriAttention} (\S\ref{sec:triattention}): Parseval 항등식에 의해 TriAttention~\cite{triattention2026}의 trig 기반 중요도 스코어는 우리의 WHT 회전 아래에서도 불변임을 관찰했다 --- 양자화된 K 블록에서 직접 계산 가능하다. 108바이트 하이브리드 블록(\texttt{AMX3\_1}) --- 50B WHT 양자화 키 + 58B polar 요약 --- 을 설계하고, 스코어링, 히스토그램 기반 Top-$B$ 임계값, 물리적 eviction을 모두 CUDA로 구현했다. 참조 구현을 지배하는 GPU$\to$CPU 왕복을 제거했고, \emph{rank 축}(누구를 남길지)과 \emph{value 축}(원소당 비트 수)을 깔끔하게 분리했다.

\item \textbf{다중 변형 엔지니어링} (\S\ref{sec:engineering}): 블록 크기 64--576에 걸친 22개 TBQ/TBQP/AMX 타입 변형을 구현했다. double WHT per-head($d=64$), 비대칭 K/V dimension(GLM 576/512), 블록별 vs.\ 전역 정규화, QJL 잔차 보정, 하이브리드 AMX 블록을 포함하며, 모두 템플릿 인스턴스화된 flash attention 프레임워크 내에서 동작한다.
\end{enumerate}

% ============================================================
\section{배경}
\label{sec:background}

\subsection{TurboQuant}

TurboQuant~\cite{turboquant2026}은 세 단계로 KV 캐시 벡터를 압축한다:

\paragraph{1. 랜덤화 WHT.} $d$차원 벡터 $\mathbf{x}$에 대해 $\mathbf{w} = H_d \cdot (\mathbf{s} \odot \mathbf{x})$를 계산한다. $H_d$는 $d \times d$ Hadamard 행렬이고 $\mathbf{s} \in \{-1, +1\}^d$는 고정 랜덤 부호 패턴이다. WHT는 입력 구조와 무관하게 계수 분포를 균등화하여, Lindeberg CLT에 의해 모든 좌표가 근사적으로 i.i.d.\ 가우시안이 된다.

\paragraph{2. Lloyd-Max 양자화.} 각 WHT 계수를 가우시안 분포에 최적인 Lloyd-Max 코드북으로 $b$비트 양자화한다. $b=3$(8레벨)의 이론적 SQNR은 31.2, $b=4$(16레벨)는 56.2이다.

\paragraph{3. QJL 보정 (선택).} TurboQuant-prod(TBQP)는 $(b{-}1)$비트 Lloyd-Max 양자화와 잔차에 대한 1비트 QJL 투영을 결합하여, 추가적인 어텐션 스코어 분산 대가로 복원 MSE를 개선한다.

\paragraph{디코딩.} 어텐션 시, 쿼리를 WHT 변환하여 키 도메인과 매칭한다: $\text{score} = \langle H_d(\mathbf{s} \odot \mathbf{q}),\ \text{dequant}(\mathbf{k}_\text{tbq}) \rangle / d$. 밸류의 경우, WHT 도메인에서 가중합을 계산한 후 역변환한다: $\mathbf{v}_\text{out} = \mathbf{s} \odot H_d(\sum_k \alpha_k \cdot \text{dequant}(\mathbf{v}_k)) / d$.

\subsection{TriAttention}

TriAttention~\cite{triattention2026}은 학습된 transformer에서 쿼리 $\mathbf{q}$가 키 $\mathbf{k}$에 할당하는 어텐션 가중치가, 단조 함수 범위 내에서, pre-RoPE $\mathbf{q}, \mathbf{k}$ 성분의 소수의 trigonometric moment로 잘 근사됨을 관찰한다. 고정 주파수 집합 $\{\omega_j\}$에 대한 레이어별·헤드별 요약 $\phi_\ell(\mathbf{k})$와 running Z-normalization을 precompute하면, 어텐션 자체를 다시 실행하지 않고도 각 토큰 $t$에 rank 스코어 $r_\ell(t)$를 부여할 수 있다. Top-$B$ 임계 미만 토큰은 KV 캐시에서 제거되어, 원소당 압축 위에 대략 $2\times$의 live-token 압축이 추가된다.

공개된 참조 구현은 이 스코어링을 호스트에서 수행하며, 트리거마다 전체 정밀도 K 캐시를 PCIe로 되읽는다. \S\ref{sec:triattention}에서 우리는 캐시가 이미 양자화된 WHT 도메인에 있을 때 rank 스코어가 우리의 부호-패턴+Hadamard 회전 아래에서 불변임을 보이며 --- 즉, 블록에서 in-place로 읽어낼 수 있다.

% ============================================================
\section{시스템 설계}
\label{sec:design}

\subsection{쿼리 전처리: GPU-Side WHT}

원본 TurboQuant 논문은 WHT를 양자화(CPU측)에서 수행한다고 가정한다. Flash attention 커널에서는 \emph{쿼리}도 런타임에 WHT 변환해야 한다. GPU 레지스터와 공유 메모리만으로 구현했다:

\paragraph{$d=256$ (128 스레드, 각 2원소).}
부호 곱셈(레지스터) → 스테이지 7(스트라이드 128, 레지스터 버터플라이) → 스테이지 0--4(\texttt{\_\_shfl\_xor\_sync} 워프 셔플) → 스테이지 5--6(공유 메모리, XOR 주소). 전역 메모리 트래픽 제로.

\paragraph{$d=512$ (128 스레드, 각 4원소).}
스테이지 8--7(스트라이드 256, 128, 레지스터) → 스테이지 0--4(워프 셔플) → 스테이지 5--6(공유 메모리). 총 9단계 버터플라이, 전역 메모리 트래픽 제로.

\subsection{K 스코어링: WHT 도메인 내적}

TBQ 키의 경우 내적을 WHT 도메인에서 직접 계산한다:
\begin{equation}
\text{score} = \sum_{i=0}^{d-1} c_{\text{idx}[i]} \cdot \text{norm} \cdot Q_{\text{wht}}[i]
\end{equation}
TBQP 키는 QJL 투영의 두 번째 항이 추가된다:
\begin{equation}
\text{score}_\text{tbqp} = \text{score}_\text{mse} + d_\text{qjl} \cdot \sum_{i} \text{sign}_\text{qjl}[i] \cdot Q_{\text{wht2}}[i]
\end{equation}

\subsection{V 누적 및 IWHT 디코딩}

밸류 벡터는 WHT 도메인에서 역양자화되고 온라인 softmax로 누적된다. KV 루프 완료 후 IWHT 역변환:
\begin{equation}
\mathbf{v}_\text{out} = \frac{1}{d}\,\mathbf{s} \odot H_d\!\left(\frac{\sum_k e^{\alpha \cdot s_k} \cdot \mathbf{v}_k}{\sum_k e^{\alpha \cdot s_k}}\right)
\end{equation}

\subsection{블록별 vs.\ 전역 정규화}

$d=512$ TBQ3의 경우 WHT 벡터의 각 256원소 절반을 독립적으로 정규화한다(\emph{블록별 norm}). 두 반쪽의 에너지가 512-point WHT 후 크게 다를 수 있기 때문이다. TBQP3는 QJL 잔차 보정이 전체 512원소 벡터에 걸쳐 동작하므로 \emph{전역 norm}을 사용한다.

% ============================================================
\section{동적 어텐션 샤프닝}
\label{sec:sharpening}

\subsection{Softmax 평탄화 문제}

K 캐시 항목의 양자화 노이즈가 각 어텐션 스코어 $s_i = \mathbf{q}^\top \mathbf{k}_i$에 독립 노이즈 $\epsilon_i$를 추가한다. 관측 스코어 $\hat{s}_i = s_i + \epsilon_i$의 softmax 분포는 $s_i$의 분포보다 평탄하다: 노이즈가 실효 온도를 높여 올바른 최댓값에 대한 확률 질량을 감소시킨다.

\subsection{$\alpha(N)$ 유도}

샤프닝된 스코어의 softmax가 원본과 일치하도록 곱셈 보정 $\alpha > 1$을 구한다.

\paragraph{MMSE 조건.} 샤프닝된 스코어의 SNR이 원본과 일치하려면:
\begin{equation}
\alpha \approx 1 + \frac{1}{2 \cdot \text{SQNR}_\text{eff}}
\end{equation}

\paragraph{EVT 보정.} $N$개 경쟁 토큰 중 최대 노이즈는 Gumbel 분포에 의해 $\sqrt{2\ln N}$으로 성장한다:
\begin{equation}
\boxed{\alpha(N) = 1 + c \cdot \sqrt{\ln N}}
\end{equation}
여기서 $c = 1/(2 \cdot \text{SQNR}_\text{eff} \cdot \sqrt{\ln N_0})$, 기준 $N_0 = 2048$.

\subsection{타입별 계수}

\begin{table}[h]
\centering
\small
\begin{tabular}{@{}lcccc@{}}
\toprule
타입 & 구조 & SQNR$_\text{eff}$ & $c$ & $\alpha$(2K) \\
\midrule
TBQ3  & 3비트 Lloyd-Max     & 31.2 & 0.00579 & 1.016 \\
TBQP3 & 2비트 + 1비트 QJL   & 13.8 & 0.01304 & 1.036 \\
TBQ4  & 4비트 Lloyd-Max     & 56.2 & 0.00326 & 1.009 \\
TBQP4 & 3비트 + 1비트 QJL   & 24.5 & 0.00724 & 1.020 \\
\bottomrule
\end{tabular}
\caption{샤프닝 계수. TBQP 타입은 QJL 랜덤 투영이 어텐션 스코어에 두 번째 노이즈 소스를 추가하므로 실효 SQNR이 낮다.}
\end{table}

% ============================================================
\section{Flash Attention의 반정밀도 병리학}
\label{sec:precision}

\subsection{버그}

업스트림 \texttt{fattn-vec} 커널은 네이티브 FP16을 지원하는 CUDA 아키텍처에서 \texttt{V\_DOT2\_F32\_F16\_AVAILABLE}을 정의한다. 활성화 시 VKQ 누적기는 \texttt{half2}(16비트), 공유 메모리 \texttt{KQ[]}는 \texttt{half}, 쿼리 레지스터 \texttt{Q\_reg}는 \texttt{half2}를 사용한다.

표준 FP16 KV 캐시에서는 이것이 정확하다. 그러나 TBQ 밸류에서는 세 가지 정밀도 손실이 복합된다:

\paragraph{손실 1: VKQ 누적.} $\sum_k \text{softmax}_k \cdot \mathbf{v}_k$가 \texttt{half2}로 누적된다. 각 덧셈이 FP16(10비트 가수부)으로 절삭되어, 수백 개 KV 항목에 걸쳐 라운딩 오차가 누적된다.

\paragraph{손실 2: 공유 메모리 스테이징.} reduce 단계에서 \texttt{\_\_shared\_\_ half KQ[]}에 부분합을 쓴 후, IWHT가 이를 \texttt{float*}로 읽는다 — 타입 불일치가 값을 추가 절삭한다.

\paragraph{손실 3: WHT 증폭.} IWHT 버터플라이는 인접 원소의 합과 차를 계산한다. $d$-point WHT에서 각 출력은 $\log_2 d$ 버터플라이 스테이지를 통해 \emph{모든} $d$개 입력에 의존한다. 한 입력의 작은 라운딩 오차 $\epsilon$이 출력에서 $O(\sqrt{d})$ 오차로 전파된다. $d=512$: $\sqrt{512} \approx 22.6$배 증폭.

\subsection{증상}

증폭된 오차가 MoE expert routing 스코어를 결정 경계 너머로 밀어, 동일 입력과 \texttt{temp=0}에서도 실행마다 다른 expert가 선택된다. 이로 인해 토큰 위치 23--27에서부터 \emph{완전히 다른 출력 시퀀스}가 생성되었다.

\subsection{수정}

모든 TBQ V 타입 템플릿 인스턴스에서 커널 헤더 전에 매크로를 undefine하여 FP32 누적을 강제한다:
\begin{verbatim}
#include "../common.cuh"
#undef V_DOT2_F32_F16_AVAILABLE
#include "../fattn-vec.cuh"
\end{verbatim}

추가로, IWHT 입력을 위한 float 레지스터 스테이징을 구현했다:
\begin{verbatim}
float tbq_regs[D_V / nthreads];
// reduce 루프가 tbq_regs에 저장
__syncthreads();
float *buf = (float*)&KQ[0];
buf[i0 + tid] = tbq_regs[i0/nthreads];
\end{verbatim}

\subsection{영향}

\begin{table}[h]
\centering
\small
\begin{tabular}{@{}lcc@{}}
\toprule
구성 & PPL & vs FP16 \\
\midrule
FP16/FP16 기준선 & 419.8 & 1.00$\times$ \\
half2 경로 (수정 전) & 454.7 & 1.08$\times$ \\
\textbf{float 경로 (수정 후)} & \textbf{445.7} & \textbf{1.06$\times$} \\
\bottomrule
\end{tabular}
\caption{FP32 누적 강제에 의한 perplexity 개선.}
\end{table}

비결정성 완전 해소: FP16 기준선은 20회 연속 동일 출력, TBQ도 이제 이 동작과 일치한다.

% ============================================================
\section{양자화된 K 위의 TriAttention}
\label{sec:triattention}

이 장에서는 TriAttention의 CUDA 네이티브 포팅과, TBQ 스타일 양자화 키와 그 스코어링에 필요한 polar 요약을 함께 담는 하이브리드 AMX3\_1 블록을 기술한다. 프레임은 이전 WHT 섹션들과 의도적으로 분리되어 있다: TBQ는 \emph{원소당 비트 수}를 줄인다(value 축); TriAttention은 \emph{보유 토큰 수}를 줄인다(rank 축). 추론 시 두 축은 곱셈적으로 합성된다.

\subsection{핵심 통찰: Parseval 불변 rank 스코어}

레이어 $\ell$에서 토큰 $t$의 TriAttention 스코어는 Z-normalization 후 내적합으로 환원된다:
\begin{equation}
r_\ell(t) = \sum_{j=1}^{J} \big( \langle \boldsymbol{\phi}^c_j, \mathbf{k}_t \rangle^2 + \langle \boldsymbol{\phi}^s_j, \mathbf{k}_t \rangle^2 \big),
\end{equation}
여기서 $\boldsymbol{\phi}^c_j, \boldsymbol{\phi}^s_j$는 주파수 $\omega_j$에 대한 고정 cos/sin 벡터. 우리의 캐시 변환 $\mathbf{k}_t \mapsto \tilde{\mathbf{k}}_t = H_d (\mathbf{s} \odot \mathbf{k}_t)$ 아래에서, Parseval 항등식에 의해
\begin{equation}
\langle \boldsymbol{\phi}, \mathbf{k}_t \rangle = \tfrac{1}{d}\langle H_d (\mathbf{s} \odot \boldsymbol{\phi}),\ \tilde{\mathbf{k}}_t \rangle,
\end{equation}
이므로 각 $\boldsymbol{\phi}$를 동일한 부호-패턴+Hadamard로 pre-rotate하면 스코어가 불변이다. 결정적으로, 회전된 $\tilde{\boldsymbol{\phi}}$는 load 시점에 한 번만 계산하면 되고, hot path는 이미 양자화된 $\tilde{\mathbf{k}}_t$와의 WHT 도메인 내적 하나만 필요하다 --- 이는 attention에 이미 계산하고 있는 내적과 정확히 일치한다. Dequantization이 불필요하다.

\subsection{하이브리드 AMX3\_1 블록 레이아웃}

평범한 TBQ3\_1 블록(128원소 WHT, 3비트 코드)은 50바이트이고 WHT 도메인 키 자체를 담기에 충분하다. 그러나 TriAttention 스코어링에는 $(r, \phi)$-form 좌표 형태의 compact \emph{polar} 요약이 추가로 필요하다 --- 전체 벡터를 재구성하지 않고 $\langle \boldsymbol{\phi}^c_j, \mathbf{k}_t \rangle, \langle \boldsymbol{\phi}^s_j, \mathbf{k}_t \rangle$를 평가하기 위함이다.

새 블록 타입 \texttt{AMX3\_1}(head\_dim $=128$, 총 108바이트)을 정의:
\begin{itemize}[leftmargin=*,topsep=2pt,itemsep=1pt]
\item \textbf{Part~A (50 B)}: 표준 TBQ3\_1 payload --- WHT 부호/Hadamard, 3비트 Lloyd-Max 인덱스, 128원소 블록당 norm 1개. Flash attention 내적에서 그대로 사용.
\item \textbf{Part~B (58 B)}: pair당 3비트 크기 $r$ + pair당 4비트 polar 각도 $\phi$ + 1비트 QJL 보정. TriAttention 스코어링 커널 전용.
\end{itemize}

Part B는 스코어링 커널(주기적 트리거, 토큰 단위 아님)에서만, Part A는 attention 커널에서만 읽는다. 두 파트는 공통 scale factor를 공유하여, Part B의 $r$이 attention 커널이 Part A에서 도출할 크기와 동일한 크기를 재현한다. 이것이 하이브리드의 내부 일관성을 만드는 핵심이다.

\subsection{CUDA 네이티브 트리거 경로}

스코어링 트리거는 매 $\Delta$ 토큰마다(기본값 256) 실행되며, 호스트 왕복 없이 캐시에 대한 3번의 GPU 패스로 진행된다:

\begin{enumerate}[leftmargin=*,topsep=2pt,itemsep=1pt]
\item \textbf{Raw 스코어 + Z-normalize.} 각 레이어와 각 주파수 $\omega_j$에 대해, 1-D 커널이 각 블록의 Part~B를 읽어 $\langle \tilde{\boldsymbol{\phi}}^c_j, \tilde{\mathbf{k}}_t \rangle^2 + \langle \tilde{\boldsymbol{\phi}}^s_j, \tilde{\mathbf{k}}_t \rangle^2$를 평가, 윈도우 running mean/variance 적용, 레이어 스코어 저장.
\item \textbf{Aggregate.} Reduce 커널이 레이어별 스코어를 토큰당 단일 rank $r(t)$로 합산.
\item \textbf{히스토그램 기반 Top-$B$ 마스크.} 장치에서 $r(t)$의 64-bin 히스토그램을 빌드; 히스토그램 prefix sum이 $|\{t : r(t) \geq \tau\}| \approx B$가 되는 임계값 $\tau$를 산출. 마스크 커널이 나머지 토큰을 eviction 대상으로 마크.
\end{enumerate}

임계값 $\tau$와 KV 위치 배열만 PCIe를 건넌다 --- 트리거당 몇 KB, 호스트 기반 참조 구현의 전체 정밀도 K 캐시(MB 단위) 대비.

\subsection{물리적 Eviction}

한 번 evicted된 블록은 다음 트리거에서 재진입되면 안 된다. WHT payload(Part~A)만 zero하는 것은 불충분하다: polar Part~B가 살아남아 다음 패스에서 블록을 $\tau$ 위로 밀어올릴 수 있다. Part~A와 Part~B(필드 \texttt{d\_wht}, \texttt{d\_r}) 모두를 zero하여 eviction을 영구화한다. 실전 차이: 없을 경우 Top-$B = 4096$에서 관찰되는 95\% eviction 비율이 수백 새 토큰 내에 70\%대로 되돌아간다.

\subsection{Attention Sink (\texttt{keep\_first})}

TriAttention만으로는 생성 안정성에 부족하다: 순수 Top-$B$ 스코어링은 결국 첫 몇 토큰을 evict하게 되고, 이때 softmax가 붕괴한다 --- 경험적으로 Qwen3-14B가 "다름, 다름, 다름\ldots" 반복 루프를 생성한다. 이는 초기 토큰이 attention sink로 작동한다는 StreamingLLM~\cite{streamingllm2023}의 관찰과 일치한다. 스코어와 무관하게 첫 $K_0$ 위치를 강제 유지하는 override 커널을 추가했다; Qwen3-14B에서 $K_0 = 32$가 충분하다(4는 너무 작고, 16은 불안정, 32는 완전 안정).

\subsection{왜 GPU 네이티브가 중요한가}

$B = 4096$에서 우리 측정은 트리거당 $\sim$2\% decode latency 비용을 보인다. 14B 규모에서 호스트 기반 참조 구현의 수 초 wall time 대비. 세 가지 요인이 기여:

\begin{itemize}[leftmargin=*,topsep=2pt,itemsep=1pt]
\item \textbf{In-place 블록 read.} 스코어링 커널은 attention에 이미 HBM에 상주하는 동일한 양자화 블록 바이트를 소비 --- dequantization 없음, PCIe 건너뜀 없음.
\item \textbf{Stream-local mask $\to$ eviction.} Mask가 GPU를 떠나지 않음 --- eviction 커널이 동일 stream에서 직접 읽음.
\item \textbf{Rank--value 직교성.} Part~A가 정확히 TBQ3\_1이므로, TriAttention 추가는 기존 flash attention 커널을 건드리지 않는다; 두 기능이 간섭하지 않고 합성된다.
\end{itemize}

\emph{메모리} 측 비용은 \S\ref{sec:results}에서 재논의: AMX3\_1은 108 B vs.\ TBQ3\_1의 50 B이므로 토큰당 원시 바이트가 증가하며, 이득은 더 적은 토큰 보유에서 온다.

% ============================================================
\section{엔지니어링 프레임워크}
\label{sec:engineering}

\subsection{타입 시스템}

블록 크기별로 조직된 20개 양자화 타입을 정의한다:

\begin{table}[h]
\centering
\small
\begin{tabular}{@{}llcl@{}}
\toprule
접미사 & 블록 크기 & Head Dim & 비고 \\
\midrule
\_0 & 256 & 256, 512 & 주력 타입 \\
\_1 & 128 & 128 & 표준 트랜스포머 \\
\_2 & 64  & 64  & 소형 head 모델 \\
\_3 & 64  & 64  & Double WHT per-head \\
\_4 & 576 & 576 & GLM 분할 256+256+64 \\
\midrule
AMX3\_1  & 128 (108\,B) & 128 & 하이브리드 WHT+polar (K측, v1.7.0) \\
AMXV3\_1 & 128 (50\,B)  & 128 & V측 AMX 동반 \\
\bottomrule
\end{tabular}
\caption{TBQ/AMX 블록 변형. 각 \_$i$ 접미사마다 TBQ3/TBQ4/TBQP3/TBQP4 하위 타입; 총 22개 변형. AMX3\_1은 TriAttention과 함께 사용되는 하이브리드 블록(\S\ref{sec:triattention}).}
\label{tab:types}
\end{table}

\subsection{Double WHT Per-Head (\_3 타입, $d=64$)}

\textbf{참고:} 이전 버전에서 기술한 double WHT per-head 접근법($H_{512} = H_8 \otimes H_{64}$)은 Q-K 도메인 불일치로 인해 \textbf{폐기}되었다. Q는 64차원 per-head 공간에서 동작하는 반면 K는 512차원 cross-head 공간에 있어 어떤 IWHT 복원도 이 간극을 해소할 수 없었다.

\paragraph{대체: Double WHT per-head.} 서로 다른 부호 패턴으로 WHT를 두 라운드 적용한다: $S_1 \to \text{WHT}_{64} \to S_2 \to \text{WHT}_{64}$. 이를 통해 첨도(kurtosis) $0.375 \to 0.047$ (거의 가우시안)을 달성하며, 단일 WHT보다 훨씬 우수한 탈상관을 보인다. Q와 K 모두 동일한 double WHT를 거치므로 Parseval 정리가 직접 적용되어 역변환이 불필요하다. 속도 영향은 무시할 수 있는 수준이다 (head당 $\sim$384 FLOP 추가).

\paragraph{$d=64$에서 QJL 필수.} 이전 분석과 달리, QJL의 1비트 보정은 멀티턴 안정성에 \emph{필수적}이다. QJL 없이(TBQ3\_3) 대화하면 3--4턴에서 반복 루프에 빠진다. QJL 포함(TBQP3\_3)시 7턴 이상 정확한 응답이 검증되었다. $d=64$에서 K는 자동으로 TBQP3\_3으로 매핑된다.

\subsection{MMSE 동적 소프트닝 ($d=64$)}

$d \geq 256$에서는 양자화 노이즈가 어텐션 gap 대비 작아 \emph{샤프닝}($\alpha > 1$)으로 softmax 평탄화를 보정한다. $d=64$에서는 노이즈가 토큰 순위 자체를 뒤집을 만큼 크다. 샤프닝은 잘못된 토큰에 대한 확신을 증폭시켜 반복 루프를 유발한다.

MMSE 최적 \emph{소프트닝} 보정을 유도했다:
\begin{equation}
\alpha(N) = \frac{\text{SQNR}}{\text{SQNR} + \sqrt{\ln N / \ln N_0}}
\end{equation}
$\alpha < 1$: softmax를 더 \emph{평탄}하게 만들어 과확신을 줄인다. EVT 항 $\sqrt{\ln N / \ln N_0}$는 긴 컨텍스트에서 더 강한 소프트닝을 적용한다.

동일한 MMSE/EVT 이론에서 높은 SQNR($d \geq 256$)에서는 샤프닝, 낮은 SQNR($d = 64$)에서는 소프트닝이 도출되는 통합 프레임워크를 형성한다.

\subsection{GLM 576: 비대칭 K/V}

GLM-4는 MLA로 $d_K=576, d_V=512$를 사용한다. \texttt{D\_V} 템플릿 파라미터로 이를 지원하며, 576차원 K는 분할 WHT(256+256+64), 512차원 V는 표준 512-point WHT를 사용한다.

\subsection{SWA F16 바이패스}

Gemma~4는 35개 레이어 중 25개에 Sliding Window Attention(SWA)을 사용한다. SWA 캐시는 300~MiB에 불과하지만(전역 KV 5120~MiB 대비), 최근 고관련성 토큰에 어텐드하므로 출력 품질을 지배한다. SWA K/V를 자동으로 FP16 업그레이드하여 가장 중요한 곳의 양자화 오차를 완전히 제거했다.

\subsection{V Rotation 비호환성}

llama.cpp KV 캐시는 기본적으로 K와 V에 RoPE rotation을 적용한다. K는 안전하다: RoPE가 직교 변환이므로 양자화기 입력의 norm과 Gaussianization 스펙트럼이 토큰 위치와 무관하게 유지되며, Parseval 항등식에 의해 상대 위치 내적 $\mathbf{q}_m^\top R_{n-m}\, \mathbf{k}$가 WHT 도메인 내적에서 보존된다(양쪽에 적용된 부호 패턴 $\mathbf{s}$는 상쇄되고, Hadamard 행렬은 직교). 그러나 V는 캐시 기록 시 적용된 RoPE를 되돌릴 역 rotation이 IWHT 디코드에 없어, 출력이 위치 의존적 회전 좌표계에 있게 된다. 모든 TBQ 구성에서 V rotation을 비활성화했다(\texttt{attn\_rot\_v=false}). 같은 논리로, 내부 WHT를 수행하는 타입(TBQ/TBQP/AMX)은 llama.cpp의 외부 \texttt{attn\_rot\_k} 경로를 반드시 우회해야 한다: 두 개의 직교 회전을 연속 적용하면 Lloyd-Max의 좌표별 centroid 가정이 깨진다.

% ============================================================
\section{개발 타임라인}
\label{sec:timeline}

구현은 8일에 걸쳐 5개 주요 세션으로 진행되었으며, 각 세션은 고유한 엔지니어링 과제를 해결했다:

\paragraph{v1.4.0 (4/1--3): 기반 + GLM 576.}
head\_dim=256(Qwen, LLaMA)과 128(표준 트랜스포머)을 위한 초기 CUDA 커널 구현. GLM-4.7-Flash의 head\_dim=576은 2의 거듭제곱이 아니어서 분할 WHT를 설계: $576 = 256 + 256 + 64$. 500회 시뮬레이션으로 full QR 직교 회전 대비 $-3.4\%$ MSE, 메모리 오버헤드 0\% 확인.

\paragraph{v1.5.0 (4/4): 업스트림 리베이스 + Gemma 4.}
업스트림 llama.cpp에 완전한 3-way merge 리베이스(26개 충돌 해결). Gemma~4: hybrid SWA(global head\_dim=512 10레이어 + SWA head\_dim=256 50레이어), 가변 GQA(16/4). SWA 캐시 타입 자동 재매핑, 가변 GQA 우회, 512-point 단일 패스 WHT(9단계 버터플라이, 128스레드$\times$4원소) 구현.

\paragraph{v1.5.1 (4/4 PM): SWA F16 바이패스.}
SWA 캐시(300~MiB)가 25레이어를 통해 출력 품질을 지배한다는 발견. SWA K/V를 FP16으로 자동 업그레이드하여 수학 벤치에서 사상 첫 f16 동급 달성(37.4/65 vs f16 36.6/65).

\paragraph{v1.5.2 세션 1 (4/7): 어텐션 샤프닝.}
PPL 21\% 격차의 원인이 비트 수가 아니라 softmax 평탄화임을 규명. 서브블록 스케일 실험 실패(WHT가 이미 분산 균등화). MMSE 이론에서 $\alpha = 1 + 1/(2 \times \text{SQNR})$ 유도. V rotation 버그 발견. PPL: 509.9 $\to$ 454.7.

\paragraph{v1.5.2 세션 2 (4/7--8): 정밀도 수정.}
temp=0 비결정적 출력 → CUDA 그래프, 풀 할당, threadfence, MoE MUL\_MAT\_ID 등 모두 원인 아님으로 제거. 근본 원인: 업스트림 \texttt{fattn-vec}의 \texttt{half2} 누적이 IWHT 버터플라이에 의해 증폭. 70개 TBQ V 템플릿 인스턴스에 FP32 강제. PPL: 454.7 $\to$ 445.7. 동적 $\alpha(N)$ 공식 구현.

\paragraph{v1.5.3 (4/11--12): $d=64$ Double WHT.}
Cross-head WHT 폐기, double WHT per-head 채택(kurtosis 0.375 $\to$ 0.047). $d=64$에서 K에 QJL 재활성화(멀티턴 안정성에 필수, 9턴+ 검증). 실험 branch \texttt{experiment/side-buffer-f16}에서 f16 side buffer로 35/35 수학 달성. $d=64$ 추천 구성: \texttt{--cache-type-k tbq4 --cache-type-v tbq3}.

\paragraph{v1.6.0 (4/14): Polar Derotate 탐구 --- 이상한 수학 결과.}
TriAttention을 알기 전, 우리는 RoPE 회전된 K를 phase와 magnitude로 명시적 분리하면 더 정확히 양자화될 것이라는 아이디어를 독자적으로 추구했다 --- \emph{polar derotate}. 프로토타입 인코더는 각 K pair를 $(r, \phi)$ 형태로 분할, $r$을 Lloyd-Max로, $\phi$를 uniform 각도 그리드로 양자화했고, dequantization 시 위치 회전을 명시적으로 되돌렸다. Qwen3-14B 수학 벤치마크(동일 35문제 suite)에서 FP16 베이스라인이 $\sim$13/35였는데, 우리의 polar-derotate 캐시는 \textbf{31/35}에 도달했다 --- TBQ3는 물론, 놀랍게도 비압축 FP16 베이스라인보다도 훨씬 높다. 이 우위를 다른 모델에서 재현하지 못했고 메커니즘도 완전히 설명하지 못한다(각도 그리드의 암묵적 정규화 효과가 의심되지만 열린 문제로 남김; \S\ref{sec:discussion} 참조). v1.7.0의 TriAttention 결과가 더 건실한 rank-축 서사를 제공한 이후 전용 polar-derotate 타입은 폐기했다; v1.7.0의 AMX3\_1 블록은 polar 요약을 다른 목적(TriAttention 스코어링)으로 재사용한다.

\paragraph{v1.7.0 (4/20--21): TriAttention 통합 + 하이브리드 AMX 블록.}
domvox의 Python TriAttention 참조 구현을 발견한 후, 전체 트리거 루프를 CUDA로 포팅(\S\ref{sec:triattention})하면서 Parseval 항등식을 사용해 양자화된 WHT 블록에서 직접 스코어를 계산했다. 새 하이브리드 블록 타입 \texttt{AMX3\_1}(108 B = 50 B TBQ3\_1 payload + 58 B polar 요약)가 양자화 키와 스코어링 커널이 필요한 통계를 모두 담는다. Qwen3-14B 262K 컨텍스트에서 유효 압축률이 $\sim$2.37$\times$ raw $\times$ $\sim$2$\times$ live-token $\approx$ 4.74$\times$에 달하며, sink 토큰 eviction을 방지하기 위해 \texttt{keep\_first} $= 32$로 설정하면 needle-in-haystack recall이 유지된다(4는 너무 작고, 16은 불안정, 32는 완전 안정). 같은 릴리즈의 housekeeping: TBQ/TBQP/AMX 타입은 외부 \texttt{attn\_rot\_k} 경로를 우회(내부에서 자체 WHT 수행; 외부 회전은 중복된 두 번째 직교 변환이었음), 폐기된 TBQX3\_1 scaffolding 제거.

% ============================================================
\section{실험 결과}
\label{sec:results}

모든 실험은 NVIDIA DGX Spark(GB10, 128GB VRAM), CUDA 12.8, Gemma~4 26B-A4B-it(MoE, UD-Q4\_K\_XL)에서 262K 컨텍스트로 수행.

\subsection{Perplexity}

\begin{table}[h]
\centering
\small
\begin{tabular}{@{}lcc@{}}
\toprule
구성 & PPL (wiki-2K) & vs FP16 \\
\midrule
FP16/FP16 & 419.8 & 1.00$\times$ \\
\midrule
v1.5.1 (샤프닝 전) & 509.9 & 1.21$\times$ \\
+ 블록별 norm & 498.2 & 1.19$\times$ \\
+ 어텐션 샤프닝 & 455.0 & 1.08$\times$ \\
+ V rotation 수정 & 448.1 & 1.07$\times$ \\
+ 동적 $\alpha$ & 454.7 & 1.08$\times$ \\
+ \textbf{FP32 누적} & \textbf{445.7} & \textbf{1.06$\times$} \\
\bottomrule
\end{tabular}
\caption{v1.5.1에서 v1.5.2까지의 점진적 PPL 개선.}
\end{table}

\subsection{수학 정확도}

35문제 수학 벤치마크(2$\times$2/3$\times$3 행렬곱, 스칼라 연산, 0--2000 토큰 필러)에서 temp=0으로 평가:

\begin{table}[h]
\centering
\small
\begin{tabular}{@{}lcccc@{}}
\toprule
구성 & 모델 & 평균 (/35) & 피크 & 10회 결과 \\
\midrule
FP16 & Gemma 4 26B & 20.1 & 21 & 19--21 \\
\textbf{TBQP3/TBQ3} & Gemma 4 26B & \textbf{19.1} & \textbf{23} & 16--23 \\
\midrule
FP16                     & Qwen3-14B & $\sim$13 & 14 & 11--14 \\
\textbf{Polar-Derotate}$^\dagger$ & Qwen3-14B & --- & \textbf{31} & 28--31 \\
AMX3\_1 + TriAttention   & Qwen3-14B & --- & 13 & 11--14 \\
\bottomrule
\end{tabular}
\caption{수학 벤치 결과. Gemma 4에서 TBQ 피크(23)가 FP16 최고(21)를 초과. Qwen3-14B에서 프로토타입 polar-derotate K 캐시($^\dagger$v1.6.0, 이후 폐기)가 \emph{설명되지 않은} 31/35 점프를 보였다 --- FP16 베이스라인 $\sim$13/35을 크게 초과. 이 효과에 대한 이론적 설명을 아직 갖지 못했고 관찰로 보고한다. 프로덕션 v1.7.0 AMX3\_1 블록은 이 벤치에서 FP16을 매칭하면서 TriAttention으로 메모리를 줄인다.}
\label{tab:math}
\end{table}

\paragraph{Needle-in-haystack recall (Qwen3-14B, v1.7.0).}
200K 토큰 컨텍스트에 6자리 비밀번호를 심고 eviction이 발생한 후 모델에게 검색을 요청한다. \texttt{keep\_first}$=4$에서 모델은 sink 토큰을 잃고 반복 루프로 퇴화; \texttt{keep\_first}$=16$에서 recall이 불안정; \texttt{keep\_first}$=32$에서 10회 시도 모두에서 안정적 recall.

\subsection{메모리}

\begin{table}[h]
\centering
\small
\begin{tabular}{@{}lrrr@{}}
\toprule
구성 & Global KV & 총 KV & 압축 \\
\midrule
FP16/FP16 & 5,120 MiB & 5,420 MiB & 1.0$\times$ \\
TBQP3/TBQ3 & 990 MiB & 1,290 MiB & 4.2$\times$ \\
\bottomrule
\end{tabular}
\caption{262K 컨텍스트에서의 KV 캐시 메모리 (Gemma 4 26B MoE).}
\label{tab:memory}
\end{table}

\paragraph{TriAttention 하의 유효 압축률.}
Qwen3-14B에서 AMX3\_1 블록은 TBQ3\_1의 50 B 대비 108 B이므로, \emph{원시} 캐시는 FP16 대비 $\sim$2.37$\times$에 불과하다 --- TBQ3 단독보다 열위. 이득은 \emph{live-token} 축에 있다: 컨텍스트가 차면 TriAttention이 슬롯의 $\sim$50\%를 evict하여 유효 $\sim$2.37$\times$ $\times$ $\sim$2$\times$ $\approx$ 4.74$\times$ end-to-end 압축을 FP16 대비 얻는다. Abstract에서 사용한 프레이밍이다. Tradeoff는 정직하다: 토큰당 바이트를 캐시당 토큰과 교환하며, rank 축이 남은 비트 수준 여유보다 더 압축되기 때문에 이긴다.

\subsection{트리거 오버헤드}

Qwen3-14B에서 $B=4096$의 TriAttention 트리거는 호출당(매 256토큰) wall-clock decode 시간의 $\sim$2\% 비용. Decode throughput은 $\sim$20.7 tok/s로 유지, TBQ3 단독과 수 퍼센트 이내. 동일 규모의 참조 호스트 기반 구현은 GPU$\to$CPU K-캐시 왕복으로 트리거당 수 초 정체.

\subsection{결정성}

FP32 누적 수정 후 temp=0으로 20회 연속 실행 시, 3비트 양자화 경계의 1--2개 보더라인 토큰을 제외하고 동일한 출력을 생성한다. FP16 기준선은 20/20 동일. 수정 전에는 토큰 23에서부터 극단적 발산이 발생했다.

% ============================================================
\section{논의: 열린 문제}
\label{sec:discussion}

\paragraph{Polar-derotate 수학 이상치.} v1.6.0 프로토타입은 FP16 베이스라인이 $\sim$13/35인 Qwen3-14B 수학 벤치마크에서 31/35에 도달했다. 비압축 모델을 초과하는 압축 캐시는 표면적으로 불가능하다. 우리가 생각해본 가능한 설명: (i) uniform 각도 그리드가 noisy attention 분포에 denoiser로 작용하는 암묵적 정규화 효과, (ii) 우리 스코어링 기준과 이 particular suite의 FP16 베이스라인 불일치, (iii) 우리의 alpha-샤프닝과 Qwen3-14B 특유 attention entropy 프로파일의 우연한 상호작용. 우리는 이 효과를 Gemma 4나 GPT-OSS에서 재현하지 못했다. 숫자를 제거하기보다 있는 그대로 관찰로 보고한다: 일어난 일이며, \emph{왜}를 이해하는 것이 숫자 자체보다 더 가치 있을 수 있다.

\paragraph{Rank 축과 Value 축의 독립성.} v1.6.0/v1.7.0 호의 교훈: 토큰 rank 압축(어떤 토큰을 남길지)과 비트 깊이 압축(각 남긴 토큰을 어떻게 저장할지)은 사실상 직교적이다. 우리의 초기 실험들은 둘을 혼동했고 --- 수학 벤치 품질을 양자화기에 비트를 추가하여 고치려 시도 --- 개선은 정체되었다. 축을 분리(value에 TBQ/AMX, rank에 TriAttention)하니 다른 축이 설계를 오염시키지 않고 각자의 한계까지 밀 수 있었다.

% ============================================================
\section{관련 연구}
\label{sec:related}

\paragraph{KV 캐시 양자화.} KIVI~\cite{kivi2024}와 KVQuant~\cite{kvquant2024}는 변환 없이 채널별/토큰별 양자화로 KV 캐시를 압축한다. TurboQuant의 WHT 전처리는 입력 구조와 무관하게 계수 분포를 균등화하여 최적 Lloyd-Max 양자화를 가능하게 한다.

\paragraph{토큰 수준 희소화.} TriAttention~\cite{triattention2026}은 우리가 활용하는 trig-요약 rank 통계를 도입했고, StreamingLLM~\cite{streamingllm2023}은 초기 토큰이 attention sink로 작동하므로 어떤 eviction 정책에서도 보호되어야 함을 확립했다. 둘 대비 본 연구의 기여는 \emph{양자화된} K 캐시와의 CUDA 네이티브 통합이다: WHT 도메인에서 rank 스코어와 원소별 양자화가 동일한 대수적 기저를 공유하므로 dequantization 없이 합성된다.

\paragraph{Flash Attention.} FlashAttention~\cite{dao2022flashattention,dao2023flashattention2}과 변형들은 타일링과 온라인 softmax로 어텐션 계산을 최적화한다. 본 구현은 단일 토큰 디코딩에 특화된 llama.cpp \texttt{fattn-vec} 커널을 확장한다.

\paragraph{어텐션 스코어 보정.} SmoothQuant~\cite{smoothquant2023} 등 outlier-aware 방법은 활성화 양자화를 다루지만 KV 양자화 노이즈에 의한 softmax 평탄화는 보정하지 않는다. 본 연구의 동적 샤프닝은 양자화 SQNR에서 유도된 최초의 원리적 보정이다.

% ============================================================
\section{결론}
\label{sec:conclusion}

TurboQuant KV 캐시 압축의 최초 프로덕션 구현을 제시하여, 최신 MoE 모델에서 4.2배 원시 메모리 절감과 6\% perplexity 비용을 달성했으며, Qwen3-14B에서 TurboQuant의 value 축 압축과 CUDA 네이티브 TriAttention의 rank 축 압축을 합성하여 추론 시 유효 $\sim$4.74$\times$ 압축으로 확장했다. 동적 어텐션 샤프닝 공식은 양자화 유도 softmax 평탄화에 대한 이론적 근거가 있는, 하이퍼파라미터 없는 보정을 제공한다. 반정밀도 증폭 병리학의 발견은 WHT 기반 방법의 수치 정밀도 요구사항이 표준 양자화와 근본적으로 다름을 강조한다. TriAttention 통합은 Parseval 항등식이 rank 수준과 value 수준 캐시 압축의 다리임을 보여준다: Lloyd-Max를 가능하게 하는 동일한 부호-Hadamard 회전이 trig-요약 스코어를 불변으로 남기므로, 두 합성은 호스트 왕복과 dequantization을 요구하지 않는다.

5개 head dimension에 걸친 22개 양자화 변형을 지원하는 전체 구현을 \url{https://github.com/AmesianX/TurboQuant}에 llama.cpp 포크로 공개한다.

% ============================================================
\section*{감사의 글}

본 연구는 NVIDIA DGX Spark에서 독립적으로 수행되었다. llama.cpp 커뮤니티와 이론적 기반을 제공한 Google DeepMind의 TurboQuant 저자들에게 감사한다. v1.7.0의 TriAttention 통합은 Mao 등~\cite{triattention2026}의 작업에 직접 기반하며, 우리의 CUDA 포팅은 오픈소스 Python 참조 구현 \texttt{domvox/triattention-ggml}에 크게 도움받았다 --- 스코어링과 eviction 루프 구조화의 선행 작업이 설계 시간을 수 주 절약해주었다.

% ============================================================
\bibliographystyle{plain}
\begin{thebibliography}{10}

\bibitem{turboquant2026}
S.~Ashkboos, A.~Mohtashami, R.~Pramanik, Y.~Li, A.~Jaggi, D.~Alistarh.
\newblock TurboQuant: Online KV Cache Quantization via Hadamard Transform.
\newblock In \emph{Proceedings of ICLR}, 2026.

\bibitem{dao2022flashattention}
T.~Dao, D.~Fu, S.~Ermon, A.~Rudra, C.~R\'{e}.
\newblock FlashAttention: Fast and Memory-Efficient Exact Attention with IO-Awareness.
\newblock In \emph{NeurIPS}, 2022.

\bibitem{dao2023flashattention2}
T.~Dao.
\newblock FlashAttention-2: Faster Attention with Better Parallelism and Work Partitioning.
\newblock In \emph{ICLR}, 2024.

\bibitem{kivi2024}
Z.~Liu, J.~Yuan, H.~Jin, S.~Zhong, Z.~Xu, V.~Braverman, B.~Chen, X.~Hu.
\newblock KIVI: A Tuning-Free Asymmetric 2bit Quantization for KV Cache.
\newblock In \emph{ICML}, 2024.

\bibitem{kvquant2024}
C.~Hooper, S.~Kim, H.~Mohammadzadeh, M.~W.~Mahoney, Y.~S.~Shao, K.~Keutzer, A.~Gholami.
\newblock KVQuant: Towards 10 Million Context Length LLM Inference with KV Cache Quantization.
\newblock \emph{arXiv preprint arXiv:2401.18079}, 2024.

\bibitem{qjl2024}
A.~Zandieh, M.~Daliri, I.~Han.
\newblock QJL: 1-Bit Quantized JL Transform for KV Cache Quantization with Zero Overhead.
\newblock In \emph{NeurIPS}, 2024.

\bibitem{smoothquant2023}
G.~Xiao, J.~Lin, M.~Seznec, H.~Wu, J.~Demouth, S.~Han.
\newblock SmoothQuant: Accurate and Efficient Post-Training Quantization for Large Language Models.
\newblock In \emph{ICML}, 2023.

\bibitem{triattention2026}
Y.~Mao et al.
\newblock TriAttention: Trigonometric Importance Scoring for Training-Free KV Cache Eviction.
\newblock \emph{arXiv preprint arXiv:2604.04921}, 2026.

\bibitem{streamingllm2023}
G.~Xiao, Y.~Tian, B.~Chen, S.~Han, M.~Lewis.
\newblock Efficient Streaming Language Models with Attention Sinks.
\newblock In \emph{ICLR}, 2024.

\end{thebibliography}

\end{document}
